\begin{tikzpicture}[font=\sffamily]
  \node[font=\sffamily\bfseries\small, ceindigo] at (6.5, 4.35)
    {DEL SALARIO NETO AL COSTE LABORAL TOTAL PARA LA EMPRESA (ESPAÑA, 2026)};

  % Tres barras comparativas horizontales ajustadas a ~13.0cm ancho total
  \node[font=\sffamily\bfseries\scriptsize, ceslate, anchor=east] at (2.8, 3.1) {Salario Neto};
  \fill[ceteal!80] (3.0, 2.75) rectangle (8.3, 3.45);
  \node[font=\sffamily\bfseries\scriptsize, white] at (5.65, 3.1) {Neto en banco ($\approx 72$--$82\%$ del Bruto)};

  \node[font=\sffamily\bfseries\scriptsize, ceslate, anchor=east] at (2.8, 1.9) {Salario Bruto};
  \fill[ceteal!55] (3.0, 1.55) rectangle (8.3, 2.25);
  \fill[ceamber!85] (8.3, 1.55) rectangle (9.7, 2.25);
  \fill[ceindigo!80] (9.7, 1.55) rectangle (10.5, 2.25);
  \node[font=\sffamily\scriptsize, cedeep] at (5.65, 1.9) {Neto trabajador};
  \node[font=\sffamily\bfseries\tiny, white] at (9.0, 1.9) {IRPF};
  \node[font=\sffamily\bfseries\tiny, white, align=center] at (10.1, 1.9) {SS\\$\approx 6,5\%$};

  \node[font=\sffamily\bfseries\scriptsize, cecrimson, anchor=east] at (2.8, 0.7) {Coste Empresa};
  \fill[cepanel] (3.0, 0.35) rectangle (10.5, 1.05);
  \draw[ceindigo, thick] (3.0, 0.35) rectangle (10.5, 1.05);
  \fill[cecrimson!85] (10.5, 0.35) rectangle (13.0, 1.05);
  \node[font=\sffamily\bfseries\scriptsize, ceindigo] at (6.75, 0.7) {Salario Bruto (100\%)};
  \node[font=\sffamily\bfseries\scriptsize, white, align=center] at (11.75, 0.7) {SS Empresa\\$\approx 30,5\%$--$32\%$};

  \node[rectangle, rounded corners=3pt, draw=ceindigo!70, fill=cecream, text width=12.4cm, align=left, inner sep=6pt, font=\sffamily\scriptsize]
    at (6.7, -0.75) {
    \textbf{Desglose de cotizaciones sociales en España (Régimen General 2026):}
    \textbf{Empresa:} Contingencias Comunes ($23,60\%$) + Desempleo indefinido ($5,50\%$) + FOGASA ($0,20\%$) + Formación Profesional ($0,60\%$) + \textbf{MEI 2026 ($0,75\%$)} + AT/EP (según tarifa CNAE, $\approx 1\%$--$3\%$). \quad
    \textbf{Trabajador:} CC ($4,70\%$) + Desempleo ($1,55\%$) + FP ($0,10\%$) + \textbf{MEI 2026 ($0,15\%$)} $=$ \textbf{6,50\%} total más retención IRPF.
  };
\end{tikzpicture}